% ---------------------------------------------------------------- security research
\section{Security research}\label{sec:research}

Kartik reports vulnerabilities to the affected organisation or through its disclosure program, and
the technical detail below goes no further than what is published at
\href{https://jkartik.in/research/}{jkartik.in/research}. The work starts by hand, with manual code
review, Burp Suite and Nmap; what is learned is then automated with Nuclei, Semgrep and BBOT.

\vspace{2pt}
{\small
\renewcommand{\arraystretch}{1.3}%
\noindent\begin{tabularx}{\linewidth}{@{}>{\raggedright\arraybackslash}p{27mm}>{\raggedright\arraybackslash}X>{\raggedright\arraybackslash}p{41mm}@{}}
\toprule[0.4pt]
\lbl{Organisation} & \lbl{Finding} & \lbl{Channel and status} \\
\midrule[0.4pt]
WhatsApp (Meta) & Privilege-escalation vulnerability in WhatsApp. & Meta bug bounty program; \$1,500 bounty awarded (2022) \\
Microsoft & Access-control and information-disclosure issues in Intune, Entra/Azure and the Exchange
  Admin Center. Organisation member data was served before the authorization check completed, a
  time-of-check/time-of-use gap spanning Intune, Entra and Azure. Standard users could bypass
  ``restricted'' policies through the Exchange admin portal to edit directory name fields and join
  closed groups. & Reported to the Microsoft Security Response Center (MSRC) \\
Mumbai Police & Stored cross-site scripting on the official website. & Disclosed responsibly, in coordination with CERT-In \\
The Souled Store & Exposed staging infrastructure that could have allowed a full system compromise. & Reported \\
Belgian Waffle Co. & Source-code and credential leak that opened the door to a potential CRM takeover. & Reported \\
Banking client & Local file inclusion and broken access control, both high severity, found through
  manual code review; detection then automated with custom Nuclei templates (\pref{exp:rudra}). & Reported to the client during a R.U.D.R.A engagement \\
\bottomrule[0.4pt]
\end{tabularx}}

\subsection{Defensive work}

\minientry{Ransomware response and infrastructure hardening}{Freelance · May -- Jun 2025}
Hardened the server infrastructure of a web platform with 16,000+ users and, when ransomware hit
it, led the response that contained the attack and restored service through rapid backup and
recovery (\pref{exp:freelance}).

\subsection{Adversarial AI}

\begin{itemize}
  \item \textbf{1st place, global leaderboard}, AI Jailbreakathon by chat.win: an LLM prompt
        attack-and-defence contest, won with the fewest tokens (\$250 prize).
  \item \textbf{Fool The LLM}: built and ran a prompt-injection contest platform (\pref{proj:ftl}).
  \item \textbf{CSAW LLM Attack Challenge (LLMAC) '25}, NYU: participated.
  \item \textbf{Writing}: \href{https://jkartik.in/blog/where-llm-guardrails-break/}{``Where LLM
        guardrails break''}, a field report from the contest on what cracks guardrail prompts and
        how to design systems that do not depend on them.
\end{itemize}
